\exercisesection{\S~6}

\begin{enumerate}
\def\labelenumi{\arabic{enumi})}
\item
  证明 \(Z(\lambda)\) 可以定义为§4习题1中的表示 \(\rho_{\lambda}\) 的一个商。
\item
  令 \(\mu\) 为 \(Z(\lambda)\)（分别地，\(E(\lambda)\)）的一个权。证明存在权序列 \(\mu_0, \ldots, \mu_n\)，属于 \(Z(\lambda)\)（分别地，\(E(\lambda)\)），使得 \(\mu_0 = \lambda, \mu_n = \mu\)，且 \(\mu_{i-1} - \mu_i \in B\) 对 \(1 \leq i \leq n\) 成立。
\item
  假设 g 是单的，并以 \(\tilde{\alpha}\) 表示 R 的最高根。模 \(\mathrm{E}(\tilde{\alpha})\) 同构于赋有伴随表示的 g。若 C 是与 g 的 Killing 型相关联的 Casimir 元，则 C 在 \(\mathrm{End}\mathrm{E}(\tilde{\alpha})\) 中的像为单位元（参见~第I章§3第7节的命题12）。利用命题7的推论推出
\end{enumerate}

\[
\varPhi_ {\mathrm{R}} (\tilde {\alpha}, \tilde {\alpha} + 2 \rho) = 1,
\]

其中 \(\Phi_{\mathrm{R}}\) 是 \(\mathfrak{h}^*\) 上的典则双线性型（第VI章§1第12节）。

\begin{enumerate}
\def\labelenumi{\arabic{enumi})}
\setcounter{enumi}{3}
\tightlist
\item
  沿用第4节的记号。
\end{enumerate}

\begin{enumerate}
\def\labelenumi{\alph{enumi})}
\item
  令 \(m \in \mathbf{N}\)。给 \(\mathbf{N}^m\) 赋予乘积序。证明对于任意 \(S\)，若它是 \(\mathbf{N}^m\) 的子集，则 \(S\) 的极小元素集是有限的。
\item
  令 \(\alpha_{1},\ldots ,\alpha_{m}\) 为 R 中互异的元素，\(X_{i}\in \mathfrak{g}^{\alpha_{i}} - \{0\}\)，令 S 为非零序列 \((p_i)\in \mathbf{N}^m\) 且满足 \(\sum p_i\alpha_i = 0\) 的集合，M 为 S 的极小元素集。则 \(\mathfrak{h}\) 与 \(X_1^{p_1}\dots X_m^{p_m}\)（其中 \((p_1,\dots ,p_m)\in \mathrm{M}\)）生成代数 \(\mathrm{U}^0\)。
\item
  证明 \(U^{0}\) 既是左诺特代数又是右诺特代数。（给 \(U^{0}\) 赋予由 \(\mathrm{U}(\mathfrak{g})\) 的滤过诱导的滤过，并利用 a) 和 b) 证明 gr \(U^{0}\) 是有限型交换代数。）
\end{enumerate}

\(d\) ) 证明对于所有 \(\lambda \in \mathfrak{h}^*\)，\(\mathrm{U}^{\lambda}\) 是有限型的左（分别地，右）\(\mathrm{U}^0\)-模。

\begin{enumerate}
\def\labelenumi{\alph{enumi})}
\setcounter{enumi}{4}
\tightlist
\item
  令 V 为满足 \(V = \bigoplus_{\lambda \in b^{*}} V_{\lambda}\) 的单 g-模。若其中某个 \(V_{\lambda} \neq 0\) 是有限维的，则所有的 \(V_{\lambda}\) 都是有限维的。（使用 d）。
\end{enumerate}

\begin{enumerate}
\def\labelenumi{\arabic{enumi})}
\setcounter{enumi}{4}
\tightlist
  \item
  证明若 \(\mathfrak{g} = \mathfrak{sl}(2,k)\)，则本段的模 \(Z(\lambda)\) 同构于§1习题2中的模 \(Z(\lambda)\)。
\end{enumerate}

